\documentclass{resume} % Use the custom resume.cls style

\usepackage{tabularx}
\usepackage{enumitem}
\usepackage{needspace}

\usepackage[left=0.4in,top=0.4in,right=0.4in,bottom=0.4in]{geometry} % Document margins
\name{Tanvir Hossain} % Your name
\address{+1(785) 424-0647 \\ Lawrence, KS (open to relocation)}
\address{\href{mailto:tanvir@ku.edu}{tanvir@ku.edu} \\ \href{https://www.linkedin.com/in/tanvir-hossain71/}{linkedin.com/in/tanvir-hossain71} \\ \href{https://scholar.google.com/citations?user=lPEOU2wAAAAJ&hl=en}{Google Scholar} \\ %
\href{https://www.tanvirhossain.net/}{tanvirhossain.net}}

\begin{document}

%----------------------------------------------------------------------------------------
%	SUMMARY
%----------------------------------------------------------------------------------------

\begin{rSection}{Summary}

Offensive hardware security researcher (Ph.D.\ expected 2027) with 4+ years of hands-on experience attacking silicon-level defenses: power/EM side-channel analysis, fault injection, hardware Trojan design, and gate-level netlist reverse engineering. Track record of published attacks and security tooling at IACR TCHES, ACM TODAES, and GLSVLSI; Champion of the IEEE HOST 2023 Microelectronics Security Challenge (IP security). Combines RTL/ASIC design-flow fluency (Verilog, Cadence and Xilinx flows) with low-level C, LLVM-based instrumentation, and machine-learning-driven attack automation. Currently building ML-guided fault-injection workflows at Keysight Technologies.
\end{rSection}

%----------------------------------------------------------------------------------------
\begin{rSection}{Technical Skills}

\begin{tabularx}{\textwidth}{@{}p{5.4cm} X@{}}
    \textsc{\textbf{Offensive Hardware Security:}} & Power/EM side-channel analysis (SCA) of cryptographic implementations, fault injection, hardware Trojan design and insertion, attacks on PUFs and obfuscated scan chains, gate-level netlist reverse engineering, threat modeling of design-for-security primitives.\\
    \textsc{\textbf{ML for Security:}} & ML-assisted side-channel analysis, golden-chip-free hardware Trojan detection and localization, model-guided fault-injection parameter optimization.\\
    \textsc{\textbf{Languages:}} & C (low-level/embedded), Python, Verilog, VHDL, Assembly (MIPS), Tcl, Bash, MATLAB; LLVM compiler-pass development.\\
    \textsc{\textbf{SCA \& Lab Equipment:}} & ChipWhisperer Husky; Keysight/Riscure Spyder, Inspector, and Pinata; high-precision EM probes and probe station; PicoScope; Tektronix DPO70404C oscilloscope; Analog Discovery 2; embedded-target bring-up and debugging.\\
    \textsc{\textbf{Digital Design \& EDA:}} & RTL design and verification; Cadence Genus, Encounter, Virtuoso; Xilinx Vivado, ISE; Intel Quartus; ModelSim; HSPICE, PSpice; gem5, Intel VTune, McPAT.\\
    \textsc{\textbf{Hardware Platforms:}} & Xilinx Artix-7 (Basys 3, Nexys A7-100T), Intel MAX10 (HaHa Board), Altera DE2-115 FPGAs; STM32, ESP32, Arduino, Raspberry Pi.\\
\end{tabularx}

\end{rSection}

%----------------------------------------------------------------------------------------
\begin{rSection}{Experience}

\begin{rSubsection}{Keysight Technologies}{May 2026 -- Aug 2026}{Machine Learning Security Intern --- Applied ML \& Fault Injection}{Santa Rosa, CA}
\item Developing model-guided optimization and adaptive-learning workflows that automatically tune fault-injection parameters, targeting higher vulnerability-detection coverage than baseline sweep-based flows.
\item Designing experiments and telemetry instrumentation; building scalable hyperparameter search and reliability metrics; delivering a prototype demonstration on production device-security test tooling.
\end{rSubsection}

\begin{rSubsection}{The University of Kansas}{Jan 2022 -- Present}{Graduate Research Assistant, Hardware Security (Advisor: Dr.\ Tamzidul Hoque)}{Lawrence, KS}
\item Designed Trojan-assisted meta-attacks that systematically locate and neutralize design-for-security primitives in microelectronic designs; demonstrated authentication bypass of ring-oscillator PUFs and unlocking of dynamically obfuscated scan chains (ACM TODAES 2026).
\item Led an in-depth security evaluation of lightweight homomorphic obfuscation schemes that protect software secrets against hardware Trojans (IACR TCHES 2026).
\item Built SPHINX, an automated framework for identifying and extracting security primitives from gate-level netlists --- reverse-engineering tooling that maps a design's defensive surface for attack planning (GLSVLSI 2026).
\item Developed golden-chip-free, ML-assisted power side-channel techniques for detecting and localizing hardware Trojans on FPGAs (Journal of Hardware and Systems Security 2026; IEEE NAECON 2023).
\item Built HOACS, a software-only mitigation that keeps secrets encoded during execution on untrusted COTS processors via residue number coding, implemented as C-library and LLVM-based program transformations for cryptographic and arithmetic workloads.
\item Proposed a gray-box, pre-deployment security-evaluation methodology for opaque COTS hardware, driven by power and EM side-channel measurements under zero-trust supply-chain assumptions.
\item Demonstrated power side-channel and fault attacks on the ASCON lightweight-cryptography standard and evaluated countermeasures; delivered hands-on power-SCA tutorials and workshops.
\item Mentored four NSF REU undergraduate researchers across two universities; co-developed and co-taught a hands-on computer hardware course using FPGAs and microcontrollers (NSF-funded).
\end{rSubsection}

\end{rSection}

%----------------------------------------------------------------------------------------
\Needspace*{10\baselineskip} % keep the Education heading with its entries (no orphaned title)
\begin{rSection}{Education}

{\bf Ph.D.\ in Electrical Engineering (Hardware Security)}, The University of Kansas \hfill {Expected 2027}\\
Dissertation research: offensive hardware security --- Trojan-assisted attacks on design-for-security primitives, side-channel analysis, and ML-assisted vulnerability detection.

{\bf M.S.\ in Electrical Engineering}, The University of Kansas \hfill {May 2025}

{\bf B.Sc.\ in Electrical and Electronic Engineering}, Ahsanullah University of Science and Technology \hfill {2016 -- 2021}

\end{rSection}

%----------------------------------------------------------------------------------------
\begin{rSection}{Selected Publications}
\newcommand{\doi}[1]{DOI: \href{https://doi.org/#1}{\nolinkurl{#1}}}

17 peer-reviewed publications (7 journal articles, 10 conference papers); full list on \href{https://scholar.google.com/citations?user=lPEOU2wAAAAJ&hl=en}{Google Scholar}.

\begin{enumerate}[itemsep=-0.15em, leftmargin=1.5em]
    \item \textbf{Tanvir Hossain}, Matthew Showers, Mahmudul Hasan, and Tamzidul Hoque, ``On the Security of Lightweight Homomorphic Obfuscation for Protecting Against Hardware Trojans,'' \emph{IACR Transactions on Cryptographic Hardware and Embedded Systems (TCHES)}, vol.\ 2026, no.\ 2, pp.\ 736--763, 2026. \doi{10.46586/tches.v2026.i2.736-763}

    \item \textbf{Tanvir Hossain}, S.\ M.\ Mojahidul Ahsan, Prabuddha Chakraborty, Swarup Bhunia, and Tamzidul Hoque, ``Attacking the Protections of Hardware using Trojan-Assisted Meta Attacks,'' \emph{ACM Transactions on Design Automation of Electronic Systems (TODAES)}, 2026. \doi{10.1145/3804452}

    \item \textbf{Tanvir Hossain}, S.\ M.\ Mojahidul Ahsan, and Tamzidul Hoque, ``SPHINX: A Framework for Security Primitive Hardware Identification and Extraction,'' \emph{Great Lakes Symposium on VLSI (GLSVLSI)}, 2026. \doi{10.1145/3787109.3815307}

    \item \textbf{Tanvir Hossain}, Ashutosh Ghimire (joint first author), Fathi Amsaad, Tamzidul Hoque, and Ahmed Sherif, ``AI-Assisted Side-Channel Analysis for Golden-Chip Free Detection and Localization of Hardware Trojans,'' \emph{Journal of Hardware and Systems Security}, 2026. \doi{10.1007/s41635-026-00182-4}

    \item Ashutosh Ghimire, Fathi Amsaad, \textbf{Tanvir Hossain}, Tamzidul Hoque, and Ahmed Sherif, ``FPGA Hardware Trojan Detection: Golden-Free Machine Learning Approach,'' \emph{IEEE National Aerospace and Electronics Conference (NAECON)}, pp.\ 181--186, 2023. \doi{10.1109/NAECON58068.2023.10365812}

    \item \textbf{Tanvir Hossain}, S.\ M.\ Mojahidul Ahsan, and Tamzidul Hoque, ``Potential and Pitfalls of Multi-valued Logic Circuits for Hardware Security,'' \emph{IEEE Dallas Circuits and Systems Conference (DCAS)}, pp.\ 1--6, 2023. \doi{10.1109/DCAS57389.2023.10130261}

    \item S.\ M.\ Mojahidul Ahsan, Muhammad Sakib Shahriar, Mrittika Chowdhury, \textbf{Tanvir Hossain}, Md Sakib Hasan, and Tamzidul Hoque, ``Accurate, Yet Scalable: A SPICE-based Design and Optimization Framework for eNVM based Analog In-memory Computing,'' \emph{IEEE/ACM International Conference on Computer-Aided Design (ICCAD)}, pp.\ 1--9, 2024. \doi{10.1145/3676536.3676827}
\end{enumerate}

\end{rSection}

%----------------------------------------------------------------------------------------
\begin{rSection}{Selected Honors \& Awards}

\begin{tabular}{lp{15.8cm}}
    \textsc{May 2023} & \href{http://www.hostsymposium.org/host_2023awards.php}{\textbf{Champion, Microelectronics Security Challenge (IP Security)}} --- IEEE International Symposium on Hardware Oriented Security and Trust (HOST 2023): first place in the IP-security track of HOST's hardware security competition. \\

    \textsc{2022, '23, '26} & \textbf{NSF Travel Grants} --- IEEE HOST Symposium (awarded three times). \\

    \textsc{January 2024} & \textbf{First Place, Poster Presentation} --- I2S Student Research Symposium, University of Kansas. \\

    \textsc{April 2023} & \textbf{David D.\ and Mildred H.\ Robb Award} --- EECS Department, University of Kansas. \\

    \textsc{March 2021} & \textbf{Best Paper Award} --- 23rd International Conference on Computer and Information Technology (ICCIT 2020). \\
\end{tabular}

\end{rSection}

%----------------------------------------------------------------------------------------

\end{document}
